\documentclass[11pt,a4paper]{article}
\usepackage[margin=22mm]{geometry}
\usepackage{amsmath,amssymb,amsthm,booktabs,hyperref}
\hypersetup{hidelinks,pdftitle={A gap in the number of changes to a tight lonely runner set}}
\newtheorem{thm}{Theorem}
\newtheorem{lem}[thm]{Lemma}
\newtheorem{cor}[thm]{Corollary}
\newcommand{\LR}{\operatorname{LR}}
\newcommand{\dist}[1]{\lVert#1\rVert}
\newcommand{\GW}{\mathrm{GW}}
\setlength{\emergencystretch}{2em}
\begin{document}
\begin{center}
{\LARGE\bfseries A gap in the number of changes\\[2pt]to a tight lonely runner set\par}
\medskip {\small Research draft \quad 26 November 2026}
\end{center}
\begin{abstract}
We consider replacing $k$ entries of the consecutive speed set by distinct integers of arbitrary size. For sufficiently large $n$, averaging over primitive rational times forces a small multiple for each deletion whenever $k\le c n/\log\log n$ and the result is tight. A previous arithmetic converse then gives the simultaneous Goddyn--Wong gcd tests. Consequently a tight modification in this range has at most $(1/\log2+o(1))\log\log n$ changes. A separate exact calculation classifies four-entry modifications for every $n\ge9$. Above the dilation at $n=9$, the first tight example has $n=148431434$ and doubles four uniquely determined speeds.
\end{abstract}
\section{Results}
For $V$ a set of $n-1$ distinct positive integer speeds let
\[
 \LR(V)=\max_{t\in\mathbb R/\mathbb Z}\min_{v\in V}\dist{vt},
 \qquad \dist{x}=\min_{j\in\mathbb Z}|x-j|.
\]
Tight means $\LR(V)=1/n$. Delete $k$ entries $r=n-s$ from $[n-1]=\{1,\ldots,n-1\}$, and insert $k$ distinct integers $v\ge n$. A pairing assigns a separate inserted multiple $v=m_s(n-s)$ to each deletion. Its Goddyn--Wong test is
\begin{equation}\label{test}
 \gcd(F,n-s)>1\qquad (F=s,s+1,\ldots,m_s s-1).
\end{equation}
Goddyn and Wong prove simultaneous sufficiency and ask about a converse with a bound on the number of changes~\cite[Sec.~3]{GW}. Zhang asks about the structure and speed bounds for arbitrary multi-entry changes~\cite[Questions 8.2--8.3]{Z}.
\begin{thm}\label{gap}
There is an absolute $c>0$ such that, for all sufficiently large $n$ and $1\le k\le c n/\log\log n$, every modification with $\LR(V)\le1/n$ has a unique divisible pairing and satisfies every test \eqref{test}. Conversely the tests imply tightness. Throughout this range,
\[
 k\le(1/\log2+o(1))\log\log n
\]
for tight configurations. In particular a non-Goddyn--Wong tight modification at sufficiently large $n$ must change more than $c n/\log\log n$ entries.
\end{thm}
The earlier range in the companion draft~\cite[Cor.~2]{C} was $k\le c(n/\log n)^{1/3}$. Our proof of the pairing below is independent of that argument. We use only the following arithmetic result of that draft: for a paired input with $2\max_s m_s s\le n$, the inequality $\LR(V)\le1/n$ forces \eqref{test} at every deletion. This dependency is part of both classification statements.
\begin{thm}\label{four}
For exactly four replacements, all tight configurations at $n\ge10$ are precisely the simultaneous accelerations satisfying \eqref{test}. None occurs at $10\le n<148431434$. At $n=148431434$ the unique configuration doubles $n-2,n-4,n-8,n-14$. At $n=9$ the only tight configuration is $2[8]$.
\end{thm}
Here $dA=\{da:a\in A\}$. At $n=7,8$ we also give an exhaustive result whenever the smallest of the four inserted speeds is $\le256$, with no assumption on the other three. At $n=7$ it is $3[6]$. At $n=8$ the possibilities are $2[7]$ and $2\{1,2,3,4,5,7,12\}$. The two sizes with all four inserted speeds above 256 remain unresolved in this draft.

\section{Counting coverage across many centres}
Fix $r=n-s\ge2$. For each integer $F\in[s,n-1]$ coprime to $r$, take $p$ modulo $r$ satisfying $Fp\equiv-1\bmod r$. The $p$ are distinct. At
\[
 t=p/r+x/(nr),\qquad x<s/F ,
\]
all baseline speeds except $r$ are at distance $>1/n$ if $x>0$ for $2r\ge n$, or $x>1/2$ otherwise. Indeed a nonzero residue other than $-1$ moves forward by $<1/r$. A speed in residue $-1$ is at most $F$ since $F+r\ge n$, and moves by $<1/r-1/n$. Other multiples of $r$ occur only when $2r<n$, and are then strictly outside the bad band for $x>1/2$. These are the single-deletion windows in~\cite[Sec.~2]{Z}. Inserted speeds must cover every one of them by their closed bad bands $\dist{vt}\le1/n$.

Use a common $x$-interval $J=(\alpha,\beta)$ of length $L$, valid at $H$ of the centres. Let $M_v$ denote the sum over those centres of the fractions of $J$ covered by the speed $v$. Necessarily $H\le\sum_{v\ {\rm inserted}}M_v$. Write $g=(v,r),q=r/g,h=v/g$. If $q>1$ and any window is met, then
\begin{equation}\label{reach}
 h\beta\ge n-q .
\end{equation}
Indeed the starting point $vp/r$ is at least $1/q>1/n$ from the next integer. Its positive displacement is $hx/(nq)$.
\begin{lem}\label{average}
For every inserted speed,
\[
 M_v\le {2r\over n}+{2r\over hL}.
\]
If $g=1$ then also $M_v\le2$. If $q=1$ then all selected centres have the same covered fraction, which is at most $2/n+2/(hL)$.
\end{lem}
\begin{proof}
Enlarge the sum to all $p\bmod r$. At fixed $x$ its integrand is $g$ times the number of $j\bmod q$ with $\dist{(j+hx/n)/q}\le1/n$. The sum has period $n/h$ in $x$ because a shift cycles the residues. Its mean over that period is $2r/n$: the $q$ translated unit intervals for $j+hx/n$ cover one full $q$-period of the band indicator. The integral on a remainder of length less than $n/h$ is at most its full-period integral $2r/h$. Divide by $L$. When $g=1$ the band contains at most two of the $r$ equally spaced points, since $2r/n<2$. For $q=1$, apply the single-period estimate to the common function $\mathbf1_{\dist{hx/n}\le1/n}$, of period $n/h$.
\end{proof}

\paragraph{Locating the deletions.}
Count the units $F$ in $I=[s,\min(n-1,\lfloor4s/3\rfloor)]$, and use $J=(1/2,3/4)$. For a nonmultiple of $r$, either $g=1$ and $h=v\ge n$, or $g\ge2$ and~\eqref{reach} implies $h\ge2n/3$ whenever any window is met. Lemma~\ref{average} gives $M_v\le14r/n$. For a multiple $v=mr$ we have $m\ge2$. In this interval $mx>1$, so contact requires $3m/4\ge n-1$. If $n\ge7$ its covered fraction is at most $2/n+6/(n-1)\le9/n$. Hence
\begin{equation}\label{mid}
 (1-9k/n)H\le14kr/n .
\end{equation}

\paragraph{Forcing small multiples.}
Assume $2r\ge n$, put $B=\lfloor n/4\rfloor\ge s$, count the units in $[s,B]$, and set $\beta=s/B$. Use the outer half $J=(\beta/2,\beta)$ of the common window. For a nonmultiple with $g=1$, Lemma~\ref{average} bounds $M_v$ by two. At $g\ge2$, contact requires $h\beta\ge n/2$, so $M_v\le10r/n\le10$. A multiple $mr$ with $m\beta>2$ has $mx>1$ throughout the window. To meet it one needs $m\beta\ge n-1$, so its covered fraction is at most $2/n+4/(m\beta)\le7/n$ for $n\ge7$. Consequently, in the absence of a multiple with $ms\le2B$,
\begin{equation}\label{outer}
 (1-7k/n)H\le10k .
\end{equation}
The change from the whole window to its outer half removes the factor $k$ from the size bound on a required multiplier.

\section{A gap in the possible numbers of changes}
Put $\rho_n=\min_{1\le r<n}\phi(r)/r$. The classical totient estimate gives $\rho_n\gg1/\log\log n$~\cite[Thm.~15]{RS}. We prove the result for $k\le n\rho_n/4096$ and sufficiently large $n$. In a block of $d$ consecutive integers there are at least $d\phi(r)/r-2^{\omega(r)}$ units modulo $r$, by inclusion--exclusion. Uniformly $2^{\omega(r)}\le 2\sqrt n=o(n\rho_n)$ by counting divisors in pairs.

If $r=1$ is deleted the interval $(1/(2n),1/n)$ is exposed. On any time interval of length $\lambda$ one speed $v$ covers at most $2\lambda/n+2/(nv)$, by period division. Using $v\ge n$ gives $n\le6k$, impossible. Now suppose $s\ge n/8$ at a deletion $r\ge2$. If $r\le s/3$, the set $I$ contains all $r$ residue classes, hence $H\ge r\rho_n$. Otherwise $r>n/24$ and $I$ has at least $n/24$ integer points: it contains $\min(r,\lfloor s/3\rfloor+1)$ points. Hence $H\ge r\rho_n/48$ for sufficiently large $n$. In both cases \eqref{mid} gives $H\le28kr/n\le28r\rho_n/4096$, a contradiction.

All offsets are now $s<n/8$. The block $[s,B]$ has at least $n/8$ entries and consequently $H\ge n\rho_n/16$. This contradicts \eqref{outer} unless some inserted multiple $m_s(n-s)$ satisfies $m_s s\le n/2$. The choices are distinct: for $a<b$ a multiple shared by $n-a,n-b$ would require a multiplier at $b$ at least $(n-a)/(b-a)>n/(2b)$. Thus the $k$ choices exhaust the inserted speeds. A speed assigned at $b$ cannot also serve any smaller offset by the same inequality. Any other bijective divisible matching would give a permutation sending each index to itself or to a larger one, hence is identical. The paired converse~\cite[Thm.~1]{C} gives all gcd tests.

For completeness these tests force $s=O(\log n)$, since every prime in $[s,2s)$ divides $n-s$, and their product has logarithm $(1+o(1))s$ by the prime number theorem. For offsets $a<b$ a prime $p\in[b,2a)$ would divide both $n-a,n-b$, hence would divide $b-a<p$. Thus $b$ exceeds the largest prime below $2a$ whenever $b<2a$. For every fixed $\epsilon>0$ successive sufficiently large offsets increase by a factor at least $2-\epsilon$. There are at most $O_\epsilon(1)+\log(O(\log n))/\log(2-\epsilon)$ offsets. Let $\epsilon\downarrow0$. This proves Theorem~\ref{gap}.

\section{Four replacements at explicit sizes}
For $k=4$, write $H_1$ for the count in \eqref{mid}, and $H_2$ for that in \eqref{outer}. A possible deletion must have $(n-36)H_1\le56r$. An upper deletion with $s\le B$ must have a multiple within budget $ms\le n/2$ if $(n-28)H_2>40n$. The distinction between $2B$ and $n/2$ can only weaken this last necessary test.

\paragraph{The infinite tail $n\ge12000$.}
Let $P_j$ denote the largest primorial $\le n$, and $\rho=\phi(P_j)/P_j$. Every $r<n$ has $\phi(r)/r\ge\rho$ and $2^{\omega(r)}\le 2^j$. Uniformly in this tail,
\[
 \rho(n-36)>56,\qquad \rho n/24-2^j>57,\qquad \rho n/8-2^j>41 .
\]
For $j=5$ insert $n=12000$ and $\rho=480/2310$. For $j=6$ insert $n=P_6=30030$ and $\rho=5760/30030$. Thereafter the lower bound $P_j\rho$ increases by a factor $p_{j+1}-1>2$ per step while $2^j$ doubles. These inequalities also give the first bound at larger $j$. If $s\ge n/8$, either $I$ contains a full residue system or $|I|\ge n/24$. In the first case $(n-36)H_1>56r$, and in the second $H_1>57>56n/(n-36)$. Thus $s<n/8$, and $H_2>41>40n/(n-28)$ forces a separate multiple within budget for each deletion. The earlier matching and converse complete the tail.

\paragraph{The finite reduction.}
For $9\le n<12000$ the file \texttt{reduce.cpp} counts units by signed squarefree divisors. At $r=1$ use $n\le24$ from the exposed interval. For $r\ge2$ discard a deletion which violates the middle inequality; mark an upper one as forced if the outer inequality supplies its budgeted multiple. An optimization for $n\ge1000$ uses $H_1<60$ for a possible deletion. For $D=\lfloor(60+2^{\omega(r)})r/\phi(r)\rfloor+1$, if $r\ge D$ and $s\ge3D$ then $I$ contains $D$ consecutive integers and too many units. If there are fewer than four possible deletions, or all are forced, the size is settled. This leaves 2802 sizes, all at most 3122.

To sharpen this test, take any list of primitive centres with a common exposed $x$-interval $(l/d,h/d)$. If four speeds cover it at all the centres, one speed covers at least one fourth of the total length counted with multiplicity. In the time variable the individual interval length is $(h-l)/(dnr)$. Period division therefore puts the qualifying speed in the finite range
\begin{equation}\label{cap}
 n\le v\le\left\lfloor {8nrd\over(n-8)(h-l)}\right\rfloor .
\end{equation}
For each candidate the exact covered length, in units of $1/(vd)$ in the $x$ coordinate, is the intersection length of
\[
 [nd(vp\bmod r)+vl,\ nd(vp\bmod r)+vh]
 \quad\text{with}\quad\bigcup_{j\in\mathbb Z}[jnrd-rd,jnrd+rd].
\]
Sum these integer lengths and multiply by four to compare with $Hv(h-l)$. At a centre met by a nonmultiple one necessarily has $v\equiv jF\bmod r$, $1\le j\le\lfloor(rd+v h)/(nd)\rfloor$. Thus testing only residues occurring for at least a quarter of the centres (overcounting repeated residues is safe) is an optional acceleration. A window deliberately skipped by the code contributes no exclusion.

In \texttt{refine.cpp}, for $n\ge75$ first discard a deletion if any tested interval bundle has no qualifying speed. Then, at an upper deletion, repeat on $s/(2B)<x<s/B$ with all budgeted multiples forbidden. If this leaves no speed, a budgeted multiple is forced. Keep at most the first 64 units producing a common window. The exact cuts, rational endpoints and one-sided skips appear in the source. This reduces the unresolved set to 169 sizes, the last being 254.

For $43\le n\le254$ use additional necessary rows of candidate speeds. For an individual exposed interval or a common bundle, its row lists speeds covering at least a quarter of its length or total length, respectively. Every solution meets every retained row. A forced deletion additionally requires a multiple within budget; an upper deletion always requires a multiple by testing $t=1/r$. The finite universe for the window rows stops at $\min(80n,80000)$; any interval requiring a larger bound by \eqref{cap} is omitted in its entirety. Divisibility rows may use one extra symbol representing all larger speeds, which relaxes the test. To seek a possible exception, mark one deletion that is lower-half or lacks a budgeted multiple. Forbid its budgeted multiples in the latter case. Enumerate choices for three other deletions, discarding a prefix when its rows cannot be hit with four allowed speeds. This is an ordinary finite hitting-set test, implemented in \texttt{screen.cpp}. At the 135 eligible sizes above 42 it leaves 1555 marked quadruples, or 1284 distinct unordered sets after forgetting the mark, all at $n\le74$.

\paragraph{Covering the remaining gaps exactly.}
For every $9\le n\le42$ we check all $\sum_{n=9}^{42}\binom{n-1}{4}=850612$ deletion sets before any filtering. For the 1284 survivors above, we also check all inserted speeds, dropping the mark. In scaled time $z=nt$ the bad bands are $[(nj-1)/v,(nj+1)/v]$. Sweep their endpoints for the baseline over $[0,n/2]$; reflection supplies the other half. Record which baseline speeds are bad on each open cell. A cell needs coverage if all its bad speeds were deleted. Since the bands are closed, covering the open cells covers their limits. Subtract the bands of each subsequently selected speed.

If $u$ insertions remain, choose a longest uncovered interval of scaled length $L$. A completing speed must cover at least $L/u$ of it and hence satisfy $v\le\lfloor2un/((n-2u)L)\rfloor$ when $n>2u$. The final speed must cover the entire connected interval and satisfies $v\le\lfloor2/L\rfloor$. If every baseline multiple of $q<n$ was removed, evaluating at $1/q$ also requires some inserted multiple of $q$. At the last step take the lcm of the outstanding requirements. The code in \texttt{cover.cpp} and \texttt{circle.hpp} uses exact endpoints throughout; it stops on empty residuals before the last step or on an exhausted arithmetic guard. Neither stop occurred. The sole covering for $9\le n\le42$ was $2[8]$ at $n=9$; no candidate above 42 in this final check covered all gaps. The dilation is tight. This establishes the structural assertion at $n\ge9$.

\paragraph{First four-fold acceleration.}
Put $N=148431434$. Every offset in \eqref{test} has $2\le s\le n/2$: otherwise it fails at 1 or on a complete residue system modulo $n-s$. The necessary modulus
\[
 d_s=\prod_{\substack{p\ {\rm prime}\\s\le p^j<2s\ \text{for some }j\ge1}}p
 \qquad\text{satisfies}\qquad n\equiv s\bmod d_s .
\]
For $64\le s\le N/2$ the certificate partitions this range into consecutive blocks $[l,h]$ with $h=\min(\lfloor5l/4\rfloor,\lfloor N/2\rfloor)$. All primes in $[h,2l)$ divide every $d_s$ in the block. Trial division supplies a product exceeding $N$ in each. For $2\le s<64$ solve the four simultaneous congruences for every four offsets by generalized CRT. The least possible size is $N$ and the only eligible offsets there are $2,4,8,14$. Direct gcd checks verify all four doublings. Their first coprime integers in increasing order from $s$ are $5,9,17,29$, so no multiplier larger than two works. The CRT minimum was already used in the earlier speed-coefficient computation. The classification gives its minimality among arbitrary four-entry modifications above the small dilation, as claimed.

\paragraph{Smaller sizes.}
For $n<5$ four changes are impossible. At $n=5$ all solutions are $d[4]$ or $d\{1,3,4,7\}$ with $d\ge5$. At $n=6$ the solutions are $d[5]$ with $3\le d\le5$ and $d\{1,3,4,5,9\}$ with $2\le d\le5$. These follow from the known classification of at most five moving speeds~\cite[Sec.~1.2]{Z}, by counting the retained entries. At $n=7,8$ fixing the \emph{smallest} new speed to each value from $n$ to 256 leaves $u=3$ and again $n>2u$. Thus the same exhaustive branch bounds cover the other three speeds without an external cutoff. Exactly the three configurations stated in Section~1 survive; each comes from a tight configuration by dilation.

There is also a relative bound in the two unresolved cases. Write the four new speeds as $M=v_1<v_2<v_3<v_4$, with $M\ge257$. Some interval in ordinary time on which all retained speeds are good has length at least $\lambda=2/21$ for $n=7$ and $\lambda=1/14$ for $n=8$. To check this, sort the endpoints $(nj\pm1)/(nv)$ of the retained bands in $[0,1]$, test their successive midpoints, and take the largest good gap for each of the 15 or 35 retained sets. After inserting the first $j$ speeds, the uncovered length on such an interval of length $\lambda$ is at least $\lambda(1-2j/n)-2j/(nM)$. It is split among at most $\lambda\sum_{i\le j}v_i+2j+1$ gaps, by counting the band centres within $1/(nv_i)$ of the interval. Use the longest gap and the $u=4-j$ remaining speeds. If $v_i\le A_i M$ and $A_1=1$, we can take
\[
 A_{j+1}=c_u\,\frac{\sum_{i\le j}A_i+(2j+1)/(257\lambda)}
 {n-2j-2j/(257\lambda)},\qquad
 c_1=2,\quad c_u=2un/(n-2u)\quad(u=2,3).
\]
The denominators are positive. The smallest remaining speed is at most the bound on a speed covering a $1/u$ share. Direct rational substitution gives $A_4<123$ at $n=7$ and $A_4<23$ at $n=8$. Thus any further example at these sizes has respectively $v_4<123M$ or $v_4<23M$.

\paragraph{Status.} Computer-assisted research draft prepared with OpenAI Codex. The finite tests are reproducible; the analytical arguments and the cited companion converse await independent review.
\begin{thebibliography}{9}\small\raggedright
\bibitem{C} \emph{A converse for sparse tight lonely runner modifications}, companion research draft, 26 November 2026. Included with the verification files.
\bibitem{GW} L. Goddyn and E. B. Wong, Tight instances of the lonely runner, \emph{Integers} \textbf{6} (2006), A38. \href{https://www.sfu.ca/~goddyn/Papers/063tight_lonely_runner.pdf}{Author's copy}.
\bibitem{RS} J. B. Rosser and L. Schoenfeld, Approximate formulas for some functions of prime numbers, \emph{Illinois J. Math.} \textbf{6} (1962), 64--94. \href{https://doi.org/10.1215/ijm/1255631807}{Published paper}.
\bibitem{Z} Y. Zhang, \emph{Tight instances of the Lonely Runner Conjecture: one-entry modifications and the growth of speeds}, \href{https://arxiv.org/html/2608.13599v3}{arXiv:2608.13599v3}.
\end{thebibliography}
\end{document}
